
\begin{abstract}

% ==== 开头段（1-2句话）====
本文针对【题目背景】问题，通过数学建模对【研究对象】进行优化设计，使得【优化目标】具有重大的现实意义。

% ==== 逐问摘要（每问按"针对问题X，我们构建了XX模型"格式）====
\textbf{针对问题一，}我们构建了【模型名称】。首先，通过【方法1】确定【中间结果1】；接着，基于【方法2】计算【中间结果2】，包括【子项1】、【子项2】以及【子项3】；然后，结合【数据】以及【公式】，代入【参数】求出了【结果类型】；最后，得到【最终结果】。最终我们得到了【关键指标1】为\textbf{【数值+单位】}，【关键指标2】为\textbf{【数值+单位】}，【关键指标3】为\textbf{【数值+单位】}，其余结果详见表X和表X。

\textbf{针对问题二，}我们构建了【模型名称】。我们以【目标函数】为目标函数，以【决策变量1】、【决策变量2】以及【决策变量3】等参数作为决策变量，以【约束条件】等要求确定了多个约束条件，建立起【模型类型】。由于决策变量的数量很多，因此我们采用【算法名称】对该模型进行求解，最终求得了【关键指标】最大为\textbf{【数值+单位】}，此时【其他结果】为\textbf{【数值+单位】}，【布局特征描述】，其余最优设计参数详见表X及resultX.xlsx。

\textbf{针对问题三，}我们构建了【模型名称】。相较于问题二，问题三中【新增约束/变化描述】，决策变量的数量进一步增多。为了简化模型，我们参考问题二中求得的结论，认为【简化假设】，又因为【性质】，所以我们可以假设【进一步简化】。因此新增的决策变量减少为【新决策变量】，其余目标函数和约束条件均与问题二相同。最终求得【关键指标】最大为\textbf{【数值+单位】}，相较于问题一的优化率为【百分比】\%，相较于问题二的优化率为【百分比】\%，其余结果详见表X及resultX.xlsx。

% ==== 结尾段（1句话）====
综上，本文为【领域】问题提供了系统的优化方案。

\keywords{关键词1；关键词2；关键词3；关键词4；关键词5；关键词6}

\end{abstract}
